% META: {"id":"TNSI-LANG-CO-003","chapitre":"TNSI-LANGAGES-ET-PROGRAMMATION","type_objet":"corrige","exercice_ref":"TNSI-LANG-EX-003","capacites_codes":["C9","C10"],"sources_inspiration":[],"status":"approved","origin":"generated","status_history":["generated","audited","accepted","approved"],"release_acceptance":"RELEASE_OWNER_BATCH_ACCEPTANCE","acceptance_closure_digest":"sha256:9b3ccf9a81c5520fbb7e03b7b2d3e7bf2057a02834b6d908bf3be5f49f3b553f","acceptance_receipt_digest":"sha256:4fad86f0282e0fa4e0419cca1ad6379ae7af5a280c04a4a90880f90a1c044242"}
% BEGIN-VERIFY
% def mots_longs_imperatif(mots, longueur_min):
%     resultat = []
%     for mot in mots:
%         if len(mot) >= longueur_min:
%             resultat.append(mot.upper())
%     return resultat
% def mots_longs_fonctionnel(mots, longueur_min):
%     return [mot.upper() for mot in mots if len(mot) >= longueur_min]
% mots = ["chat", "ordinateur", "cle", "programmation"]
% assert mots_longs_imperatif(mots, 6) == mots_longs_fonctionnel(mots, 6)
% assert mots_longs_fonctionnel(mots, 6) == ["ORDINATEUR", "PROGRAMMATION"]
% END-VERIFY
\begin{corrige}{TNSI-LANG-EX-003}
\begin{python}
def mots_longs_fonctionnel(mots, longueur_min):
    return [mot.upper() for mot in mots if len(mot) >= longueur_min]
\end{python}
\begin{python}
mots = ["chat", "ordinateur", "cle", "programmation"]
print(mots_longs_imperatif(mots, 6) == mots_longs_fonctionnel(mots, 6))  # True
print(mots_longs_fonctionnel(mots, 6))  # ['ORDINATEUR', 'PROGRAMMATION']
\end{python}
Les deux versions produisent le même résultat : la version fonctionnelle exprime en une seule expression ce que la version impérative construit pas à pas avec une variable d'accumulation (\lstinline{resultat}) modifiée par des instructions successives.
\end{corrige}
